%-----------------------------------------------------------------------------------------------------
\section{Lesson 1, 23/09}
\textbf{PRIMA LEZIONE SU TABLET}
\section{Lesson 2, 24/09}
\subsection{Natural units}
From this point onwards, natural units will be used:
$$\hbar = c = 1$$
\textbf{Note:} this messes up dimensionality, but, convenientely, forces mass and energy to have
the same units
$$E=mc^2\underset{c=1}\rightarrow E=m,$$
and lenght and time to have the same dimensionality
$$x=ct\underset{c=1}\rightarrow[x]=[t].$$
To convert back into more familiar units the constants are, respectively,
$$c=299.792.458\text{ m/s}$$
$$\hbar c \approx 0.97 \text{ GeV$\cdot$fm}$$
\subsection{Klein-Gordon Equation}
The relativistic quantum mechanics' fundamental equation is
$$(\Box+m^2)\psi(x)=0$$
where
$$\Box = \partial^\mu\partial_\mu=\pdv[2]{}{t}-\laplacian$$
and $\psi(x)$ is still the Schrödinger's wavefunction.

It will now be shown that this equation \textbf{does not} conserve probability, using a continuity
equation.

Regular quantum mechanics states
$$\pdv[]{\rho}{t}+\div\mathbf{j} = 0,$$
where
$$\rho=\bra{\psi_t}\ket{\psi_t}\quad\quad \mathbf{j=}\frac{i}{m}\left(\psi^*\grad\psi-\psi\grad\psi^*
\right)$$

Consider the Klein-Gordon equation left-multiplied by $\psi^*(x)$ and its conjugate
$$\begin{dcases}
  \psi^*(x)\left(\Box+m^2\right)\psi(x)=0\\
  \psi(x)\left(\Box+m^2\right)\psi^*(x)=0
  \end{dcases}$$
their difference is
$$\psi^*(\Box+m^2)\psi - \psi(\Box+m^2)\psi^*=0$$
$$\partial^\mu(\psi^*\partial_\mu\psi - \psi\partial_\mu\psi^*)=0$$
$$\partial^\mu J_\mu=0$$
this is a 4D continuity equation and $J^\mu$ is a Noether current.

Expanding the equation yields
$$\pdv[]{}{t}\left(\psi^*\pdv[]{\psi}{t}-\psi\pdv[]{\psi^*}{t}\right)-\div\left(\psi^*\grad \psi
-\psi\grad\psi^*\right)=0$$
the second term is proportional to the quantum probability current, but it is missing an $i$!
Moreover, the first term, representing $\rho$, is \textbf{not definite positive}, so it cannot
represent the norm of the state $\ket{\psi_t}$!\\

This proves that it is \textbf{impossible} to make QM relativistic using wavefunctions and retaining
the probabilistic approach.

Paul Dirac used another approach, he wanted to decrease the order of space derivative (instead of
increasing the time derivative order, like K-G): this result was better, but for free particles
it yielded a double solution with $E<0$ states, violating a fundamental QM theorem!\\
For this reason, he devised the concept of the \textit{Dirac sea}:
these $E<0$ states must all be occupied (because they cannot be observed), and a uniform charge
can be factored out of Maxwell's equations with proper boundary conditions; when one of these
particles in the sea is excited, the "hole" left behind can be percieved as a "positive electron":
he predicted the existence of \textit{antiparticles}.

This interpretation is still classical, but it begins to convey the concept that a finite number
of particles is \textit{not} possible for relativistic QM, so the concept of wavefunctions must be
left behind in favor of field theories.
\subsection{Symmetries and conservations}
To build a theory, it is necessary to start with principles and symmetry:
\textbf{Lorentz-invariance} to guarantee the conservation of distance between points (this is also
useful to define a metric topology!).

Why symmetry? Noether theorem states that for every 4-current that respects a continuity equation
$$\partial_\mu J^\mu=0$$
there is a conserved Noether charge
$$Q:=\int_{}^{}\dd^3xJ^0\text{, so that }\dv[]{Q}{t}=0$$
\subsubsection{Lorentz-Poincaré group}
The Lorentz invariant is
$$s^2=t^2-x_ix^i = t'^2-x_i'x'^i = x_\mu x^\mu:$$
the Lorentz group is the set of all transformations $x^\mu\rightarrow x'^\mu$ over the Minkowski
space that conserve $s^2$; in matricial representation the elements are
$$x'^\mu = \Lambda\indices{^\mu_\nu} x^\nu$$
the invariant is conserved \textbf{if and only if}
$$\boxed{g_{\mu\nu}\Lambda\indices{^\mu_\rho}\Lambda\indices{^\nu_\sigma} = g_{\rho\sigma}.}$$
The covariant quantities transform with the inverse matrix
$$x'_\mu = \Lambda\indices{_\nu^\mu} x'_\nu$$
\textbf{Observation 1} Since $g_{\mu\nu}$ is a symmetric matrix, this equation encompasses $10$
constraints on $\Lambda\indices{^\mu_\nu}$, leaving $3$ degrees of freedom for rotations and $3$ for
boosts.

\textbf{Naming convention} Consider the determinant of the invariance condition:
$$\det(\Lambda^Tg\Lambda)=\det(g)\leftrightarrow\det(\Lambda)^2=1\leftrightarrow\det(\Lambda)=\pm1$$
Transformations with $\det(\Lambda) = 1$ are called \textit{proper Lorentz transformations}.\\
Transformations with $\det(\Lambda)=-1$ are called \textit{imporper Lorentz transformations}.

\textbf{Naming convention} Consider the (0,0) component of the invariance condition:
$$g_{\mu\nu}{\Lambda^\mu}_0{\Lambda^\nu}_0 = 1 \leftrightarrow (\Lambda_{00})^2 - \Lambda\indices{^i_0}
g_{ii}{\Lambda^i}_0=1\leftrightarrow |{\Lambda^0}_{0}|^2\geq1$$
Because this term regulates how the transformation relates $x^0$ and $x'^0$,
transformations with $\Lambda\indices{^0_0}\geq1$ are called
\textit{orthocronous Lorentz transformations}, while those with $\Lambda\indices{^0_0}\leq1$ are
called \textit{non-orthcoronous Lorentz transformations}.

These 2 divisions split the Lorentz group into 4 parts
\begin{itemize}
    \item Proper and orthocronous $L^{\uparrow}_{+}$: restricted Lorentz group, sometimes referred to
	as "Lorentz Group";
    \item Improper and orthocronous $L^{\uparrow}_-$: spacial inversion;
    \item Proper and non-orthocronous $L^{\downarrow}_+$ : time inversion;
    \item Improper and non-orthocronous $L^{\downarrow}_-$: total inversion of spacetime.
\end{itemize}
\textbf{Note} The only subgroup is $L^\uparrow_+$, as it is the only one that contains the identity
matrix $\mathbb{I}$.
